\chapter{Research Results}
\label{ch:rezultaty}

\section{Control System Integration and Plant Definition}
\label{sec:integracja_systemu}

The architecture of the research environment constitutes a synthesis of the elements developed in the previous parts of this work. According to the general closed-loop control scheme (Figure \ref{fig:najwazniejszy_graf}), the system consists of the nonlinear dynamic model of the parallel robot derived in Chapter \ref{ch:model_dynamiczny} and the DeePC controller whose structure is defined in Chapter \ref{ch:deepc}.

\subsection{System Input and Output Definition}

From the perspective of the DeePC regulator, the system is defined solely by the available input and output signals. The control input vector $\boldsymbol{u}$ of dimension $m=2$ corresponds to the torques generated by the actuators in the active joints:
\begin{equation}
    \boldsymbol{u} =
    \begin{bmatrix}
        \tau_1 \\ \tau_4
    \end{bmatrix} \in \mathbb{R}^2
    \label{eq:input_definition}
\end{equation}
The measurement output vector $\boldsymbol{y}$ of dimension $p=2$ contains only the angular positions of the active arms:
\begin{equation}
    \boldsymbol{y} =
    \begin{bmatrix}
        \theta_1 \\ \theta_4
    \end{bmatrix} \in \mathbb{R}^2
    \label{eq:output_definition}
\end{equation}
This deliberately limited observability — excluding velocity measurements and direct end-effector localization — constitutes a more demanding test case, forcing the controller to reconstruct system behavior solely from the input-output history contained in the Hankel matrix.

\section{Control Quality Assessment Methodology}

To quantitatively evaluate control quality, the following indicators were defined \cite{ethzurich, choosewisely}:
\begin{itemize}
    \item \textbf{Mean Absolute Error} ($p$):
    \begin{equation}
        \label{eq:precision}
        p = \frac{1}{N_{sim}} \sum_{k=0}^{N_{sim}} |\boldsymbol{y}_k - \boldsymbol{y}_{ref,k}|
    \end{equation}

    \item \textbf{Maximum Control Error} ($p_{max}$):
    \begin{equation}
        \label{eq:max_precision}
        p_{max} = \max_{k} \left( |\boldsymbol{y}_k - \boldsymbol{y}_{ref,k}| \right)
    \end{equation}

    \item \textbf{Total Control Cost} ($c_{total}$):
    \begin{equation}
        c_{total} = \sum_{k=0}^{N_{sim}} \left( \| \boldsymbol{y}_k - \boldsymbol{y}_{ref,k} \|_{\boldsymbol{Q}}^2 + \| \boldsymbol{u}_k \|_{\boldsymbol{R}}^2 \right)
    \end{equation}

    \item \textbf{Average Optimization Time} ($t_{solve}$) and \textbf{Total Computational Time} ($t_{sel}$) for Selective DeePC:
    \begin{equation}
        t_{sel} = t_{solve} + \frac{1}{N_{sim}} \sum_{k=0}^{N_{sim}} t_{sort,k}
    \end{equation}
\end{itemize}

The indicators $p$ and $p_{max}$ were defined analogously for the task space as $p_{pos}$ and $p_{pos,max}$, computed as the Euclidean distance between the reference and actual end-effector coordinates obtained from forward kinematics.

\section{Research Environment and Implementation}
\label{sec:srodowisko}

All simulations were conducted in MATLAB on a platform equipped with a 12th Gen Intel Core i7-12650H processor and 32 GB of RAM. The YALMIP library \cite{yalmip} and the MOSEK solver \cite{mosek} were used to solve the optimization problem.

The base framework was taken from \cite{zhang_deepc_code, github1}. The most important original contributions include:
\begin{itemize}
    \item \textbf{Implementation of Selective DeePC} using the methodology of \cite{choosewisely}.
    \item \textbf{Mosaic Hankel matrix construction} from multiple independent experiments via (\ref{eq:mosaic_hankel}).
    \item \textbf{Backtrack mechanism} for safe data generation within workspace boundaries.
    \item \textbf{Adaptation to the parallel robot}, including singular configuration avoidance constraints in the QP formulation.
\end{itemize}

\section{Data Preparation and Optimization Formulation}
\label{sec:generacja_danych}

The offline historical database was generated using a random walk method adapted from \cite{choosewisely}. Fifty independent experiments of 100 steps each were conducted, yielding a total dataset of $T=5000$ samples combined into a mosaic Hankel matrix via (\ref{eq:mosaic_hankel}). The exciting signal followed:
\begin{equation}
    \boldsymbol{u}_k = a \cdot \boldsymbol{u}_{k-1} + b \cdot \boldsymbol{n}_k
\end{equation}
with parameters $a=1$, $b=0.1$, $n_a=0.4$, where $\boldsymbol{n}_k$ is drawn from a uniform distribution on $[-n_a, n_a]$. A backtrack mechanism ($n_{back}=10$) ensured all collected trajectories remain kinematically feasible. The resulting point cloud is shown in Figure \ref{fig:A_CHMURA}.

\begin{figure}[!htbp]
    \centering
    \includegraphics[width=0.9\linewidth]{rysunki/4.rezultaty/Chmura_punktów.pdf}
    \caption{Visualization of the training data point cloud against the reference trajectory in the robot workspace}
    \label{fig:A_CHMURA}
\end{figure}

The physical actuator and geometric constraints were defined as:
\begin{align}
    u_{i} &\in [-0.1,\ 0.1]\,\mathrm{Nm}, \quad i=1,2 \\
    y_{1} &\in \left[\tfrac{\pi}{8},\ \tfrac{11}{24}\pi\right], \quad
    y_{2} \in \left[-\tfrac{11}{24}\pi,\ -\tfrac{\pi}{8}\right]
\end{align}
All signals were normalized to $[-1, 1]$ via:
\begin{equation}
    \label{eq:normalization}
    x_{n} = 2\frac{x - x_{min}}{x_{max} - x_{min}} - 1
\end{equation}
Hard box constraints were imposed on both inputs and outputs in the QP task:
\begin{align}
    \boldsymbol{H}_u \boldsymbol{u}_k &\leq \boldsymbol{h}_u, \quad
    \boldsymbol{H}_y \boldsymbol{y}_k \leq \boldsymbol{h}_y
\end{align}
with $\boldsymbol{H}_u = [\boldsymbol{I}; -\boldsymbol{I}]_{4\times2}$ and $\boldsymbol{h}_u = \boldsymbol{1}_{4\times1}$, analogously for outputs. A rate constraint $|\Delta \boldsymbol{u}_k| \leq \Delta u_{max} = 0.4$ was additionally imposed to limit control aggressiveness.

\section{Experimental Configuration}
\label{sec:konfiguracja_eksperymentow}

Controller parameters were tuned on Trajectory A (Lemniscate of Gerono — a smooth $C^\infty$ curve) and subsequently evaluated on Trajectory B. The common simulation parameters were:
\begin{equation}
    T_s = 0.01\,\mathrm{s}, \quad T_{sim}=5\,\mathrm{s}, \quad T_{ini}=8, \quad N=20
\end{equation}
\begin{equation}
    \boldsymbol{Q} = 300 \cdot \boldsymbol{I}_{p}, \quad \boldsymbol{R} = 0.01 \cdot \boldsymbol{I}_{m}
\end{equation}

\textbf{Trajectory B (Polygon)} was the primary evaluation scenario (Figure \ref{fig:robotrownoleglyB}, Table \ref{tab:wspolrzedne_B}). Its vertices are traversed at a constant time of $1\,\mathrm{s}$ per segment ($100$ steps at $T_s=0.01\,\mathrm{s}$), imposing variable velocities and velocity discontinuities at corners. This physically unrealizable reference was intentionally designed to stress-test algorithm robustness under extreme dynamics.

\begin{figure}[!htbp]
    \centering
    \includegraphics[width=1\linewidth]{rysunki/4.rezultaty/druga_trajektoria.pdf}
    \caption{Kinematic scheme of the robot with active link constraints and Trajectory B (Polygon)}
    \label{fig:robotrownoleglyB}
\end{figure}

\begin{table}[!htbp]
    \centering
    \caption{Vertex coordinates of reference Trajectory B}
    \label{tab:wspolrzedne_B}
    \begin{tabular}{ccc}\toprule
        \textbf{Vertex No.} & \textbf{Coordinate} $x$ [m] & \textbf{Coordinate} $y$ [m]\\\midrule
        0 (Start) & $x_0$ & $y_0$ \\
        1 & $0.06$ & $0.00$ \\
        2 & $0.10$ & $0.01$ \\
        3 & $0.20$ & $-0.01$ \\
        4 & $0.16$ & $-0.06$ \\
        5 & $0.06$ & $0.00$ \\\bottomrule
    \end{tabular}
\end{table}

\section{Results -- Trajectory B, No Measurement Noise}

Controller parameters for the noise-free case:
\begin{itemize}
    \item \textbf{Regularized DeePC:} $\lambda_g=100,\ \lambda_{\sigma}=10^8$
    \item \textbf{Selective DeePC:} $\lambda_g=0.1,\ \lambda_{\sigma}=10^8,\ N_{cols}=155$
\end{itemize}

\subsection{Workspace Trajectory Analysis}

The tracking results in the XY plane are presented in Figure \ref{fig:bez_szumu_porownanie_b}. The corner cutting visible at vertices is not an algorithm error but a physically necessary behavior: since the reference requires an instantaneous velocity direction change, the DeePC controller generates a realizable trajectory at the cost of a temporary geometric deviation.

\begin{figure}[!htbp]
    \centering
    \includegraphics[width=0.9\linewidth]{rysunki/4.rezultaty/BEZ SZUMU/Trajektoria_B_bez_szumu.pdf}
    \caption{Trajectory B tracking in the workspace (no noise)}
    \label{fig:bez_szumu_porownanie_b}
\end{figure}

The Selective DeePC method shows a clear advantage in corner precision — the local linearization mechanism enables faster adaptation to direction changes. The Regularized controller exhibits overshoots at vertices, a consequence of the averaging effect introduced by the full historical database. In the second-to-last segment (approx. 3--4 s), both methods show local quality degradation, consistent with a data shortage for this specific motion region near singular configurations.

\subsection{Quantitative Assessment}

\begin{table}[!htbp]
    \centering
    \caption{Joint space control quality indicators -- Trajectory B (no noise)}
    \label{tab:bledy_katy_B}
    \renewcommand{\arraystretch}{1.2}
    \begin{tabular}{c w{c}{2.8cm} w{c}{2.8cm} w{c}{2.8cm} w{c}{2.8cm}}
        \toprule
        \multirow{2}{*}{\textbf{Method}} &
        \multicolumn{2}{c}{\textbf{Tracking Precision $p$} [$\cdot 10^{-3}$ rad]} &
        \multicolumn{2}{c}{\textbf{Maximum Error $p_{max}$} [$\cdot 10^{-3}$ rad]} \\
        \cmidrule(lr){2-3} \cmidrule(lr){4-5}
        & $\theta_1$ & $\theta_4$ & $\theta_1$ & $\theta_4$ \\
        \midrule
        Regularized & $1.96$ & $2.39$ & $8.58$ & $8.20$ \\
        Selective   & $\mathbf{1.77}$ & $\mathbf{1.28}$ & $\mathbf{5.80}$ & $\mathbf{7.07}$ \\
        \bottomrule
    \end{tabular}
\end{table}

\begin{table}[!htbp]
    \centering
    \caption{End-effector positioning error indicators -- Trajectory B (no noise)}
    \label{tab:bledy_pos_B}
    \renewcommand{\arraystretch}{1.2}
    \begin{tabular}{c c c}
        \toprule
        \textbf{Method} &
        \textbf{Mean Error $p_{pos}$} [mm] &
        \textbf{Maximum Error $p_{pos,max}$} [mm] \\
        \midrule
        Regularized & $0.54$ & $2.08$ \\
        Selective   & $\mathbf{0.42}$ & $\mathbf{2.03}$ \\
        \bottomrule
    \end{tabular}
\end{table}

The Selective method dominates across all indicators. The mean angular error for $\theta_4$ is nearly twice as small as in the Regularized method, and the task-space precision is superior in both mean and maximum error (Table \ref{tab:bledy_pos_B}). The Euclidean error profile (Figure \ref{fig:bez_szumu_blad_euklidesowy_b}) confirms that outside the problematic 3--4 s segment, Selective maintains a consistently lower error level.

\begin{figure}[!htbp]
    \centering
    \includegraphics[width=0.9\linewidth]{rysunki/4.rezultaty/BEZ SZUMU/Porownanie_Blad_Euklidesowy_B_bez_szumu.pdf}
    \caption{End-effector Euclidean error over time -- Trajectory B (no noise)}
    \label{fig:bez_szumu_blad_euklidesowy_b}
\end{figure}

\subsection{Control Signal Analysis}

The control signal comparison (Figure \ref{fig:bez_szumu_sygnal_b}) reveals the fundamental operational difference between the two methods. Selective DeePC generates signals with larger amplitude changes, enabling faster cornering response at the cost of chattering. The Regularized method produces smoother profiles, tempered by the global data matrix, at the expense of geometric fidelity.

\begin{figure}[!htbp]
    \centering
    \includegraphics[width=0.9\linewidth]{rysunki/4.rezultaty/BEZ SZUMU/Porownanie_wejscia_B_bez_szumu.pdf}
    \caption{Comparison of control signals for Trajectory B (no noise)}
    \label{fig:bez_szumu_sygnal_b}
\end{figure}

\section{Results -- Trajectory B, With Measurement Noise}
\label{sec:nastawy_szum}

Additive white Gaussian noise was introduced to the joint angle measurements to simulate real encoder imperfections:
\begin{equation}
    \tilde{y}_{i,k} = y_{i,k} + \xi_{i,k}, \quad \xi_{i,k} \sim \mathcal{N}(0, \sigma_i^2)
\end{equation}
with standard deviation $\sigma_i = 0.005 \cdot (y_{max,i} - y_{min,i})$, corresponding to $0.5\%$ of the full operating range. Controller parameters were retuned accordingly:
\begin{itemize}
    \item \textbf{Regularized DeePC:} $\lambda_g=20000,\ \lambda_{\sigma}=10^8$
    \item \textbf{Selective DeePC:} $\lambda_g=200,\ \lambda_{\sigma}=5\cdot10^5,\ N_{cols}=155$
\end{itemize}

\subsection{Workspace Trajectory Analysis}

The noisy trajectory comparison (Figure \ref{fig:z_szumem_porownanie_b}) reveals a fundamental divergence in algorithm behavior:

\begin{figure}[!htbp]
    \centering
    \includegraphics[width=0.9\linewidth]{rysunki/4.rezultaty/Z SZUMEM/Trajektoria_B_szum.pdf}
    \caption{Comparison of polygon shape reproduction in the presence of noise}
    \label{fig:z_szumem_porownanie_b}
\end{figure}

\begin{itemize}
    \item \textbf{Regularized DeePC:} The corner cutting phenomenon is drastically intensified. The robot fails to reach the trajectory turnaround points, tracing curvilinear segments inscribed within the polygon. The bias is significant and variable, resulting from the algorithm averaging noisy data across the full historical database.
    \item \textbf{Selective DeePC:} Despite the noise, this method maintains significantly higher geometric fidelity. Through the first three segments the robot follows the straight lines closely and approaches the vertices. Stability problems appear only in the 4th and 5th segments.
\end{itemize}

\subsection{Quantitative Assessment}

\begin{table}[!htbp]
    \centering
    \caption{Joint space control quality indicators -- Trajectory B (with noise)}
    \label{tab:bledy_katy_B_szum}
    \renewcommand{\arraystretch}{1.2}
    \begin{tabular}{c w{c}{2.8cm} w{c}{2.8cm} w{c}{2.8cm} w{c}{2.8cm}}
        \toprule
        \multirow{2}{*}{\textbf{Method}} &
        \multicolumn{2}{c}{\textbf{Tracking Precision $p$} [$\cdot 10^{-2}$ rad]} &
        \multicolumn{2}{c}{\textbf{Maximum Error $p_{max}$} [$\cdot 10^{-2}$ rad]} \\
        \cmidrule(lr){2-3} \cmidrule(lr){4-5}
        & $\theta_1$ & $\theta_4$ & $\theta_1$ & $\theta_4$ \\
        \midrule
        Regularized & $2.75$ & $2.36$ & $\mathbf{5.59}$ & $7.52$ \\
        Selective   & $\mathbf{2.23}$ & $\mathbf{1.72}$ & $6.95$ & $\mathbf{6.98}$ \\
        \bottomrule
    \end{tabular}
\end{table}

\begin{table}[!htbp]
    \centering
    \caption{End-effector positioning error indicators -- Trajectory B (with noise)}
    \label{tab:bledy_pos_B_szum}
    \renewcommand{\arraystretch}{1.2}
    \begin{tabular}{c c c}
        \toprule
        \textbf{Method} &
        \textbf{Mean Error $p_{pos}$} [mm] &
        \textbf{Maximum Error $p_{pos,max}$} [mm] \\
        \midrule
        Regularized & $\mathbf{6.15}$ & $\mathbf{17.61}$ \\
        Selective   & $6.85$ & $46.04$ \\
        \bottomrule
    \end{tabular}
\end{table}

The mean Cartesian error of Selective is inflated by a local instability in the 4th segment, where oscillations introduce a large initial bias at the start of the 5th segment, visible as the $46.04\,\mathrm{mm}$ peak in Figure \ref{fig:z_szumem_blad_euklidesowy_b}. Outside this critical region, Selective maintains an error below $5\,\mathrm{mm}$ for most of the trajectory, compared to a $10\,\mathrm{mm}$ baseline for Regularized. This highlights the trade-off characteristic of the selective approach: superior mean performance with a risk of local destabilization in data-sparse regions.

\begin{figure}[!htbp]
    \centering
    \includegraphics[width=0.9\linewidth]{rysunki/4.rezultaty/Z SZUMEM/Porownanie_Blad_Euklidesowy_B_szum.pdf}
    \caption{End-effector Euclidean error -- Trajectory B (with noise)}
    \label{fig:z_szumem_blad_euklidesowy_b}
\end{figure}

\subsection{Control Signal Analysis}

The control signal profiles (Figure \ref{fig:z_szumem_sygnal_b}) confirm the contrasting strategies: Selective DeePC performs continuous small corrections to maintain geometric fidelity, while Regularized DeePC generates smooth but geometrically inaccurate outputs. This behavior is consistent across both noise conditions and constitutes the defining operational characteristic of each method.

\begin{figure}[!htbp]
   \centering
   \includegraphics[width=0.9\linewidth]{rysunki/4.rezultaty/Z SZUMEM/Porownanie_wyjsc_B_szum.pdf}
   \caption{Control signal profiles -- Trajectory B (with noise)}
   \label{fig:z_szumem_sygnal_b}
\end{figure}

\section{Computational Performance}

Computation times are summarized in Tables \ref{tab:czasy_bez_szumu} and \ref{tab:czasy_szum}. The results are consistent across both noise conditions, confirming that noise presence does not affect computational complexity.

\begin{table}[!htbp]
    \centering
    \caption{Average computation times -- Trajectory B (no noise)}
    \label{tab:czasy_bez_szumu}
    \renewcommand{\arraystretch}{1.4}
    \begin{tabular}{c c}
        \toprule
        \textbf{Method} & \textbf{Trajectory B} \\
        \midrule
        Regularized DeePC & $t_{solve} = 108.38~\mathrm{ms}$ \\
        Selective DeePC & $t_{sel} = \mathbf{14.65}~\mathrm{ms}$ \quad ($t_{sort} = 1.57~\mathrm{ms}$) \\
        \bottomrule
    \end{tabular}
\end{table}

\begin{table}[!htbp]
    \centering
    \caption{Average computation times -- Trajectory B (with noise)}
    \label{tab:czasy_szum}
    \renewcommand{\arraystretch}{1.4}
    \begin{tabular}{c c}
        \toprule
        \textbf{Method} & \textbf{Trajectory B} \\
        \midrule
        Regularized DeePC & $t_{solve} = 99.96~\mathrm{ms}$ \\
        Selective DeePC & $t_{sel} = \mathbf{13.64}~\mathrm{ms}$ \quad ($t_{sort} = 1.55~\mathrm{ms}$) \\
        \bottomrule
    \end{tabular}
\end{table}

The Selective method is approximately 7 times faster in both conditions. The overhead from the data selection stage ($t_{sort} \approx 1.6\,\mathrm{ms}$) is negligible compared to the gain from reducing the decision variable dimension from $\approx3650$ to $N_{cols}=155$. These results position Selective DeePC as the more promising candidate for real-time implementation, with computation times approaching the $50\,\mathrm{Hz}$ control regime.

\section{Hyperparameter Sensitivity Analysis}
\label{sec:wrazliwosc}

Sensitivity analysis was performed on noisy Trajectory B as the most demanding test case, using the nominal parameter set from Section \ref{sec:nastawy_szum} as the reference.

\subsection{Impact of Regularization Coefficient $\lambda_g$}

\begin{figure}[!htbp]
    \centering
    \includegraphics[width=0.9\linewidth]{rysunki/4.rezultaty/WPŁYW ZMIANY PARAMETRÓW/Analiza_wpływu_lambda_g.pdf}
    \caption{Impact of regularization parameter $\lambda_g$ on control errors}
    \label{fig:wplyw_lambda_g}
\end{figure}

The analysis (Figure \ref{fig:wplyw_lambda_g}) reveals clearly distinct stability thresholds for the two methods:
\begin{itemize}
    \item \textbf{Selective DeePC} requires $\lambda_g > 10^2$ for stable operation. Above $10^3$, a gradual increase in control cost is observed, indicating a well-defined optimal range.
    \item \textbf{Regularized DeePC} requires significantly stronger regularization at $\lambda_g > 10^4$ — a threshold two orders of magnitude higher than Selective.
\end{itemize}
This difference reflects the fundamental distinction between the two approaches: by operating on a local, pre-filtered subset of data, Selective DeePC inherently requires less regularization to suppress noise, as the data selection step already acts as an implicit filter.

\subsection{Impact of Selected Column Number $N_{cols}$ (Selective DeePC)}

\begin{figure}[!htbp]
    \centering
    \includegraphics[width=0.9\linewidth]{rysunki/4.rezultaty/WPŁYW ZMIANY PARAMETRÓW/Analiza_wpływu_Ncols_Selective.pdf}
    \caption{Impact of $N_{cols}$ on Selective DeePC performance}
    \label{fig:wplyw_ncols}
\end{figure}

The analysis (Figure \ref{fig:wplyw_ncols}) reveals three distinct operating regimes:
\begin{itemize}
    \item \textbf{Precision optimum:} The lowest cost function value is recorded at $N_{cols}=170$, where the local model best approximates the robot dynamics from the most relevant historical data.
    \item \textbf{Gradual degradation:} In the range $170$--$275$ columns, a steady increase in control cost is observed as progressively less representative trajectories are included.
    \item \textbf{Instability threshold:} Beyond $N_{cols}=275$, a rapid loss of stability occurs with errors reaching hundreds of millimeters, confirming that including distant, poorly matched data is more harmful than operating on a smaller but well-selected sample.
\end{itemize}
Computation time increases monotonically with $N_{cols}$, consistent with the growing QP problem dimension, reinforcing the importance of selecting the smallest $N_{cols}$ that preserves control quality.

\section{Discussion}

The simulation studies confirm a clear performance differentiation between the two DeePC variants on the nonlinear parallel robot benchmark.

Regularized DeePC exhibits a non-zero systematic bias directly attributable to the global data approach: since the five-bar linkage dynamics are highly nonlinear, the full Hankel matrix cannot simultaneously represent all workspace regions with equal fidelity. The slack variable $\boldsymbol{\sigma}$ partially compensates for this, but at the cost of tracking accuracy. This limitation is significantly amplified under noise, manifesting as the characteristic polygon-to-curve degradation observed in Figure \ref{fig:z_szumem_porownanie_b}.

Selective DeePC addresses this limitation through local dynamics linearization. The improvement in geometric fidelity — even under noise — confirms the hypothesis that restricting the data matrix to nearest-neighbor trajectories reduces the approximation error of the implicit LTI assumption. The cost of this strategy is a higher sensitivity to data-sparse regions: when the local model is poorly conditioned, temporary instability can occur, as observed in the 4th trajectory segment.

A crucial practical finding is the sevenfold reduction in computation time achieved by Selective DeePC, from approximately $100\,\mathrm{ms}$ to $14\,\mathrm{ms}$. This reduction — driven purely by the smaller QP dimension — brings real-time implementation within reach without sacrificing, and in most cases improving, tracking quality.

The sensitivity analysis confirms that parameters tuned on Trajectory A transfer successfully to the more demanding Trajectory B, suggesting that the optimal hyperparameter configuration is primarily governed by robot dynamics rather than specific trajectory excitation. The $\lambda_g$ threshold analysis further underscores the structural advantage of the selective approach: its implicit noise filtering through data selection relaxes the regularization requirements by two orders of magnitude relative to the global method.